\subsection{模型三的建立与求解：偏差检验驱动的选择性融合}\label{sec:model-q3}

\subsubsection{模型三的建立}\label{sec:q3-establishment}

\noindent\textbf{（1）实际观测的时间配准。}

设机器人在参考时间轴上的真实轨迹为 \(\mathbf p^\star(t)\)。以方式 1
的时钟为参考，问题三的两组实际定位观测满足

\begin{equation}
\begin{aligned}
\mathbf z_i^{(1)}
&=\mathbf p^\star\!\left(t_i^{(1)}\right)
+\boldsymbol\varepsilon_i^{(1)},\\
\mathbf z_j^{(2)}
&=\mathbf p^\star\!\left(t_j^{(2)}-\Delta t\right)
+\mathbf b(t_j^{(2)})
+\boldsymbol\varepsilon_j^{(2)},
\end{aligned}
\label{eq:q3-observation}
\end{equation}

其中，\(\Delta t\) 为固定时间偏差，\(\mathbf b(t)\) 为待判定的相对系统偏差。
方式 2 的设备时刻映射到参考时间轴为

\begin{equation}
t_{2,\mathrm{ref}}=t_2-\Delta t.
\label{eq:q3-time-map}
\end{equation}

由于两组设备时间区间不直接重叠，先定义各自的经过时间

\begin{equation}
s_i^{(m)}=t_i^{(m)}-t_1^{(m)},\qquad m=1,2,
\label{eq:q3-elapsed-time}
\end{equation}

再对中心化二维位置特征计算归一化互相关，得到粗时间偏差 \(h_c\)。
以 \(h_c\) 为中心构造精搜索域 \(\mathcal D\)，并沿用
式\eqref{eq:q2-fixed-grid}--\eqref{eq:q2-calibration-estimator}
的固定公共网格和 Huber 剖面目标，得到

\begin{equation}
\widehat{\Delta t}
=\arg\min_{\delta\in\mathcal D}J_H(\delta).
\label{eq:q3-time-estimator}
\end{equation}

时间对齐后，在公共评价网格 \(\{t_k\}_{k=1}^{N}\) 上构造未作空间校正的
相对残差

\begin{equation}
\mathbf d_k=
\widetilde{\mathbf z}_2(t_k+\widehat{\Delta t})
-\widetilde{\mathbf z}_1(t_k).
\label{eq:q3-aligned-residual}
\end{equation}

没有外部真值时，\(\mathbf d_k\) 只能用于辨识两设备之间的相对空间差，
不能判断误差应归因于哪一设备。

\noindent\textbf{（2）HAC-Wald 偏差检验。}

检验假设为

\begin{equation}
H_0:\mathbb E(\mathbf d_k)=\mathbf 0,
\qquad
H_1:\mathbb E(\mathbf d_k)\ne\mathbf 0.
\label{eq:q3-hypothesis}
\end{equation}

记样本均值及滞后 \(r\) 阶样本自协方差为

\begin{equation}
\overline{\mathbf d}
=\frac{1}{N}\sum_{k=1}^{N}\mathbf d_k,
\qquad
\widehat{\boldsymbol\Gamma}_r
=\frac{1}{N}\sum_{k=r+1}^{N}
(\mathbf d_k-\overline{\mathbf d})
(\mathbf d_{k-r}-\overline{\mathbf d})^\mathsf T.
\label{eq:q3-autocovariance}
\end{equation}

采用 Newey--West 核构造样本均值的 HAC 协方差估计
\cite{neweywest1987}：

\begin{equation}
\widehat{\boldsymbol\Sigma}_{\overline d}
=\frac{1}{N}
\left[
\widehat{\boldsymbol\Gamma}_0+
\sum_{r=1}^{L}
\left(1-\frac{r}{L+1}\right)
\left(
\widehat{\boldsymbol\Gamma}_r+
\widehat{\boldsymbol\Gamma}_r^\mathsf T
\right)
\right],
\label{eq:q3-hac-covariance}
\end{equation}

其中 \(L\) 为 HAC 截断阶数。Wald 统计量为

\begin{equation}
W=\overline{\mathbf d}^{\,\mathsf T}
\widehat{\boldsymbol\Sigma}_{\overline d}^{-1}
\overline{\mathbf d},
\qquad
W\mathop{\sim}_{H_0}\chi_2^2.
\label{eq:q3-wald}
\end{equation}

\noindent\textbf{（3）工程效应与移动块 Bootstrap。}

统计显著不必然意味着工程上需要校正。令

\begin{equation}
r_k=\|\mathbf d_k\|_2,\qquad
s_r=1.4826\,
\operatorname{median}_k
\left|r_k-\operatorname{median}_{\ell}r_\ell\right|,
\label{eq:q3-residual-scale}
\end{equation}

定义工程效应指数

\begin{equation}
E_b=\frac{\|\overline{\mathbf d}\|_2}{s_r}.
\label{eq:q3-engineering-effect}
\end{equation}

设统计显著性水平为 \(\alpha\)，工程效应阈值为 \(\eta\)，偏差启用判据为

\begin{equation}
I_b=
\mathbf 1\!\left\{
p_{\mathrm{HAC}}<\alpha
\ \land\
E_b\ge\eta
\right\}.
\label{eq:q3-bias-decision}
\end{equation}

为避免依赖渐近分布，同时保留残差的局部时序结构，采用移动块
Bootstrap\cite{kunsch1989}。给定块长 \(\ell\)，构造首尾相接的残差块

\begin{equation}
\mathcal B_j=
\left(
\mathbf d_j,\mathbf d_{j+1},\ldots,\mathbf d_{j+\ell-1}
\right),
\label{eq:q3-bootstrap-block}
\end{equation}

每次有放回抽取若干完整块并截取前 \(N\) 个残差，计算重抽样均值
\(\overline{\mathbf d}^{*(b)}\)。重复 \(B\) 次后，分别取两个坐标方向的
\(\alpha/2\) 与 \(1-\alpha/2\) 分位数，得到逐轴置信区间

\begin{equation}
\mathrm{CI}_r=
\left[
Q_{\alpha/2}
\!\left(\overline d_r^{\,*}\right),
Q_{1-\alpha/2}
\!\left(\overline d_r^{\,*}\right)
\right],
\qquad r\in\{x,y\}.
\label{eq:q3-bootstrap-ci-estimator}
\end{equation}

\noindent\textbf{（4）偏差趋势与状态维度选择。}

为了区分固定偏差与缓慢漂移，对每个坐标拟合

\begin{equation}
d_{kr}=\beta_{0r}+\beta_{1r}t_k+u_{kr},
\qquad r\in\{x,y\},
\label{eq:q3-trend-model}
\end{equation}

并使用 HAC 协方差检验 \(H_0:\beta_{1r}=0\)。令 \(I_{\mathrm{tr}}\)
表示至少一个坐标方向检出显著趋势，则状态模型按

\begin{equation}
\mathcal M=
\begin{cases}
\mathcal M_0,& I_b=0,\\
\mathcal M_{\mathrm c},& I_b=1,\ I_{\mathrm{tr}}=0,\\
\mathcal M_{\mathrm{rw}},& I_b=1,\ I_{\mathrm{tr}}=1
\end{cases}
\label{eq:q3-model-selection}
\end{equation}

选择。其中，\(\mathcal M_0\) 为无偏差状态，
\(\mathcal M_{\mathrm c}\) 为常量偏差状态，
\(\mathcal M_{\mathrm{rw}}\) 为随机游走偏差状态。

无偏差状态与问题二一致：

\begin{equation}
\mathbf x_k^{(0)}
=\left(
\mathbf p_k^\mathsf T,
\mathbf v_k^\mathsf T,
\mathbf a_k^\mathsf T
\right)^\mathsf T.
\label{eq:q3-state-no-bias}
\end{equation}

启用偏差时扩展为

\begin{equation}
\mathbf x_k^{(b)}
=\left(
\mathbf p_k^\mathsf T,
\mathbf v_k^\mathsf T,
\mathbf a_k^\mathsf T,
\mathbf b_k^\mathsf T
\right)^\mathsf T.
\label{eq:q3-state-with-bias}
\end{equation}

扩展状态的测量矩阵为

\begin{equation}
\mathbf H_1=
\begin{bmatrix}
\mathbf I_2&\mathbf 0&\mathbf 0&\mathbf 0
\end{bmatrix},
\qquad
\mathbf H_2=
\begin{bmatrix}
\mathbf I_2&\mathbf 0&\mathbf 0&\mathbf I_2
\end{bmatrix}.
\label{eq:q3-bias-observation-matrix}
\end{equation}

常量偏差模型满足
\(\mathbf b_k=\mathbf b_{k-1}\)；随机游走模型满足

\begin{equation}
\mathbf b_k=\mathbf b_{k-1}+\boldsymbol\xi_k,
\qquad
\boldsymbol\xi_k\sim
\mathcal N(\mathbf 0,q_b\tau_k\mathbf I_2).
\label{eq:q3-bias-random-walk}
\end{equation}

\noindent\textbf{（5）选择性异步融合。}

两设备测量协方差仍由式\eqref{eq:q2-third-difference}和
式\eqref{eq:q2-measurement-covariance}估计。将方式 2 的事件时刻按
式\eqref{eq:q3-time-map}校正后，与方式 1 的原始事件合并排序。
根据式\eqref{eq:q3-model-selection}选择状态维度，再执行
式\eqref{eq:q2-kalman-prediction}--\eqref{eq:q2-rts}的异步 Kalman
滤波和 RTS 平滑。

过程噪声调参阶段关闭门控反馈。对白化创新
\(\mathbf e_k=\operatorname{chol}(\mathbf S_{0,k})^{-1}\boldsymbol\nu_k\)，
在候选集 \(\mathcal Q\) 上最小化

\begin{equation}
S(q)=
\left|\overline{\gamma}(q)-2\right|
+2\,\overline{\left|\rho_1(\mathbf e;q)\right|},
\qquad
\widehat q=\arg\min_{q\in\mathcal Q}S(q).
\label{eq:q3-process-score}
\end{equation}

固定 \(\widehat q\) 后运行 NIS 稳健门控。对于 10 Hz 非事件时刻，从前一
RTS 平滑事件状态按状态方程传播。由输出协方差的位置子矩阵
\(\mathbf P_{p,k}\) 定义二维合成位置标准差

\begin{equation}
\sigma_{p,k}=
\sqrt{\operatorname{tr}(\mathbf P_{p,k})}.
\label{eq:q3-position-std}
\end{equation}

\subsubsection{模型三的求解}\label{sec:q3-solution}

\begin{pseudocodeblock}{算法3\quad 偏差检验驱动的选择性融合}
{\heiti 输入：}两路实际定位序列
\(\mathcal Z_1,\mathcal Z_2\)，粗搜索集 \(\mathcal H\)，
精搜索半宽 \(L_t\)，显著性水平 \(\alpha\)，工程阈值 \(\eta\)，
HAC 阶数 \(L\)，Bootstrap 次数 \(B\) 与块长 \(\ell\)，
过程噪声候选集 \(\mathcal Q\)。\par
{\heiti 输出：}时间偏差 \(\widehat{\Delta t}\)、偏差判定
\(I_b\)、所选状态模型 \(\mathcal M\) 及 10 Hz 融合状态。\par
\begin{enumerate}[label={\heiti 步骤~\arabic*:},leftmargin=5.5em,itemsep=1pt,topsep=3pt]
  \item 按式\eqref{eq:q3-elapsed-time}重置两路时间，利用中心化位置特征的
  归一化互相关获得粗时间偏差 \(h_c\)；
  \item 在 \([h_c-L_t,h_c+L_t]\) 内最小化固定网格 Huber 剖面目标，
  由式\eqref{eq:q3-time-estimator}得到 \(\widehat{\Delta t}\)；
  \item 按式\eqref{eq:q3-aligned-residual}构造同步残差，计算
  \(\overline{\mathbf d}\)、HAC 协方差、Wald 统计量和工程效应指数；
  \item 重复 \(B\) 次移动块 Bootstrap，按
  式\eqref{eq:q3-bootstrap-ci-estimator}构造逐轴置信区间；
  \item 拟合式\eqref{eq:q3-trend-model}并执行 HAC 趋势检验，按
  式\eqref{eq:q3-bias-decision}和式\eqref{eq:q3-model-selection}
  选择状态模型；
  \item 由位置三阶差分估计两设备测量协方差，关闭门控反馈并按
  式\eqref{eq:q3-process-score}选择 \(\widehat q\)；
  \item 按校正事件时刻执行所选模型的异步 Kalman 滤波与 NIS 稳健门控，
  再执行 RTS 后向平滑；
  \item 在公共支持域内按 0.1 s 间隔传播状态，输出位置、速度、加速度及
  式\eqref{eq:q3-position-std}定义的模型内位置标准差。
\end{enumerate}
\end{pseudocodeblock}

时间配准、偏差检验、状态模型选择和事件融合之间的关系如
图\ref{fig:q3-framework}所示。
\begin{figure}[htbp]
  \centering
  \includegraphics[width=.98\textwidth]{q3/method_flow.pdf}
  \caption{问题三偏差检验与选择性融合流程}
  \label{fig:q3-framework}
\end{figure}

\subsubsection{结果分析}\label{sec:q3-results}

\noindent\textbf{（1）时间配准结果。}

两种定位方式分别包含 1381 个和 1621 个有效观测。经过时间上的中心化位置
互相关得到粗时间偏差

\begin{equation}
h_c=-368.82\ \mathrm s.
\label{eq:q3-coarse-result}
\end{equation}

以 \(h_c\) 为中心取 10 s 精搜索半宽，搜索区间为
\([-378.82,-358.82]\,\mathrm s\)。固定公共评价网格含 2910 个时刻，
Huber 剖面连续精修得到

\begin{equation}
\widehat{\Delta t}=-367.877619\ \mathrm s,
\qquad
t_{2,\mathrm{ref}}
=t_2-(-367.877619)
=t_2+367.877619\ \mathrm s.
\label{eq:q3-time-offset}
\end{equation}

因此，方式 2 的设备时钟相对于方式 1 参考时钟落后约
\(367.88\,\mathrm s\)，映射时需将方式 2 时间戳整体增加相应数值。

固定网格 Huber 剖面的伴随位置中心为

\begin{equation}
\widehat{\mathbf b}_{\mathrm{prof}}
=(-0.151775,-0.150330)^\mathsf T\ \mathrm m.
\label{eq:q3-profile-bias-result}
\end{equation}

相应标定 RMSE 为 \(6.36335\,\mathrm m\)。该位置中心仅作为时间剖面目标中的
伴随参数，系统偏差是否启用由后续统计检验决定。

\noindent\textbf{（2）系统偏差统计判定。}

时间配准后用于偏差推断的残差序列含 3000 个 10 Hz 时刻，其普通均值为

\begin{equation}
\overline{\mathbf d}
=(-0.133731,-0.157069)^\mathsf T\ \mathrm m,
\qquad
\|\overline{\mathbf d}\|_2=0.206288\ \mathrm m.
\label{eq:q3-mean-residual}
\end{equation}

HAC 截断阶数取 \(L=8\)，得到样本均值协方差

\begin{equation}
\widehat{\boldsymbol\Sigma}_{\overline d}
=
\begin{bmatrix}
0.020304&0.001455\\
0.001455&0.018041
\end{bmatrix}
\mathrm m^2.
\label{eq:q3-hac-result}
\end{equation}

取 \(\alpha=0.05\)、工程阈值 \(\eta=0.25\)，检验结果见
表\ref{tab:q3-decision}。Wald 检验未拒绝零相对偏差，工程效应指数也未达到
启用偏差状态的阈值。

\begin{table}[htbp]
\centering
\caption{问题三相对系统偏差的统计与工程判定}
\label{tab:q3-decision}
\begin{tabular}{lcc}
\toprule
判定量 & 计算结果 & 判定阈值\\
\midrule
HAC-Wald 统计量 \(W\) & \(2.0935\) & \(\chi^2_{2,0.95}\)\\
\(p_{\mathrm{HAC}}\) & \(0.3511\) & \(0.05\)\\
稳健噪声尺度 \(s_r\) / m & \(3.0300\) & --\\
工程效应指数 \(E_b\) & \(0.0681\) & \(0.25\)\\
\bottomrule
\end{tabular}
\end{table}

移动块 Bootstrap 取 \(B=2000\)、块长 \(\ell=14\)，所得 95\%
置信区间为

\begin{equation}
b_x\in[-0.413289,0.138470]\ \mathrm m,
\qquad
b_y\in[-0.428052,0.119683]\ \mathrm m.
\label{eq:q3-bootstrap-ci}
\end{equation}

两轴区间均覆盖零。图\ref{fig:q3-bias-ci}同时给出了样本均值和移动块
Bootstrap 区间，零水平线均穿过两个坐标方向的置信区间。

\begin{figure}[htbp]
  \centering
  \includegraphics[width=.76\textwidth]{q3/bias_confidence_interval.pdf}
  \caption{问题三相对偏差估计及 95\% 移动块 Bootstrap 置信区间}
  \label{fig:q3-bias-ci}
\end{figure}

x、y 方向的 HAC 趋势斜率分别为
\(-1.4675\times10^{-3}\,\mathrm{m/s}\) 和
\(1.4894\times10^{-4}\,\mathrm{m/s}\)，对应 \(p\) 值为
0.3783 和 0.9208，均未检出显著漂移。图\ref{fig:q3-bias-series}
显示残差存在明显的局部时序相关，因此偏差均值的协方差不能按独立样本估计；
两条滚动均值曲线则未呈现贯穿观测区间的单调漂移。

\begin{figure}[htbp]
  \centering
  \includegraphics[width=.86\textwidth]{q3/bias_time_series.pdf}
  \caption{问题三时间配准后的两坐标相对残差时序}
  \label{fig:q3-bias-series}
\end{figure}

综合 \(p_{\mathrm{HAC}}\)、\(E_b\)、Bootstrap 区间和趋势检验，在给定
显著性水平与工程阈值下，未发现具有统计与工程意义的相对系统偏差，
即 \(I_b=0\)，选择六维无偏差状态模型 \(\mathcal M_0\)。
该结论表示现有样本不足以支持启用偏差状态，不等同于证明真实偏差严格为零。

\noindent\textbf{（3）噪声参数与异步融合。}

位置三阶差分的 MAD 估计给出

\begin{equation}
\widehat{\mathbf R}_1
=\operatorname{diag}(16.0645,18.2067)\ \mathrm{m}^2,
\qquad
\widehat{\mathbf R}_2
=\operatorname{diag}(7.2399,8.9392)\ \mathrm{m}^2.
\label{eq:q3-measurement-noise-result}
\end{equation}

过程噪声候选集取 \([10^{-6},10^3]\) 上的对数网格。
式\eqref{eq:q3-process-score}在

\begin{equation}
\widehat q=3.1623
\label{eq:q3-process-noise-result}
\end{equation}

处取得最小值，该值位于候选区间内部。图\ref{fig:q3-tuning}给出了
过程噪声谱密度与 NIS/白化创新综合得分的关系。

\begin{figure}[htbp]
  \centering
  \includegraphics[width=.80\textwidth]{q3/process_noise_tuning.pdf}
  \caption{问题三基于门控前创新统计的过程噪声选择}
  \label{fig:q3-tuning}
\end{figure}

固定 \(\widehat q\) 后采用 99\% NIS 门控。门控前平均 NIS 为
2.0127、95\% 分位为 5.9590，门控后平均 NIS 为 1.9957；
白化创新平均绝对一阶相关为 0.0274，共有 1.033\% 的观测触发协方差放大。

\noindent\textbf{（4）10 Hz 融合状态与不确定性。}

融合轨迹的公共输出区间为
\(445.0076\sim814.0076\,\mathrm s\)，按 0.1 s 间隔共生成
3691 个状态点。位置、速度、加速度和协方差派生量均为有限值；
10 Hz 非事件时刻从前一平滑事件节点传播的最大时间间隔为 0.208 s。

二维合成位置标准差的中位数为 \(0.9606\,\mathrm m\)，最大值为
\(3.4593\,\mathrm m\)。图\ref{fig:q3-trajectory}以颜色表示
\(\sigma_{p,k}\)，并标记公共输出区间的起点和终点。该标准差只反映所建
状态空间模型内的不确定性，不等同于融合轨迹相对于未知真实轨迹的绝对误差。

\begin{figure}[H]
  \centering
  \includegraphics[width=.72\textwidth]{q3/fused_trajectory_uncertainty.pdf}
  \caption{问题三 10 Hz 融合轨迹及模型内位置标准差}
  \label{fig:q3-trajectory}
\end{figure}
